\section{The separate rank-one construction: a polynomial averaging mechanism}
\label{sec:type-one}

\status{The Type-I block lattices, their primitive-test interpretation, the
actual gamma degree and leading-form identities, logarithmic integrality,
and the polynomial-matrix averaging lemma are proved below.
The full rank-one local theorem remains conditional on the restricted
graded Gram, simultaneous-partner and paired-lift inputs specified at the end.
The earlier $r\geq2$ segment theorem is not used for this branch.}

\paragraph{TI1--TI2: the actual polynomials and integral lattices.}
The edge rank is one, but its segment polynomials have $r=m$ coordinates,
where the ambient dimension is $m+1$.  A tie joins adjacent vector and
covector rows.  With scalar variables $Z_i,S_i$ and an auxiliary row $\xi$,
the source uses
\[
 E_i=Z_iS_i+\sum_c x_{a_i,c}x_{a_i+1,c},\qquad
 E_+=Z_+S_++\sum_c x_{u,c}\xi_c,
\]
and the integral span of
\[
 \Bigl(\prod_iE_i^{j_i}Z_i^{z_i}\Bigr)E_+^{j_+}Z_+^{z_+}f(x)g(\xi).
 \tag{TI1.1}
\]
Here the levels fit both tied degrees, $f$ belongs to the actual integral
prefix space of residual degrees, and $\deg g=B-j_+$.
The block lattice is its rational span intersected with
\[
 v_p(F_\alpha)+\tfrac12v_p(\alpha!)+\tfrac12v_p(s_b)\geq0,
 \qquad s_b=\prod_{i<\sigma}(A_i-w_i)!\,\mathrm{sg}_i!.
 \tag{TI1.2}
\]
This is a full finite free lattice: it is a rational subspace intersected
with a finite free weighted-monomial lattice.  Its form $s_b\,\mathrm{Fock}$
is integral term by term.  Smith bases show that its largest saturated
subspace in which every primitive vector has a partner of valuation at
most $D$ has rank exactly $\#\{i:v_p(d_i)\leq D\}$.
Indeed the indicated Smith vectors work; a larger saturated subspace
contains a primitive vector with all those coordinates zero.

\paragraph{TI3: the contraction polynomial.}
Let a signed chain have $N+1$ rows of degrees $d_i$, total $S$, and last
vector row $w$ of degree $o$. Fix $\rho$, a nonnegative integer shift $z$,
and append a covector $y$; all row degrees are nonnegative integers.
The source substitution $\mathcal F$ is
\[
\begin{array}{lll}
\text{vector }x:&x_\rho\mapsto\sum_cx_cy_c,&x_c\mapsto y_\rho x_c\ (c\ne\rho),\\
\text{covector }x:&x_\rho\mapsto y_\rho x_\rho,
 &x_c\mapsto y_\rho x_c-x_\rho y_c\ (c\ne\rho).
\end{array}
\]
Put $\operatorname{fr}_n(f)=y_\rho^{n-S}\mathcal F(f)$ for $n\geq S$,
$E=\sum_cw_cy_c$, $F=\sum_c\partial_{w_c}\partial_{y_c}$ and
$A=o+n+r+z-1$.  With $B$ denoting the bilinear Fock form, set
\[
 \Phi_{o,n,z}=\sum_{j=0}^o\frac{(-1)^jE^jF^j}{j!^2\binom Aj},\qquad
 \Gamma_{f,g}(n)=
 \frac{\binom Ao B(\Phi_{o,n,z}\operatorname{fr}_n f,
                         \operatorname{fr}_n g)}
 {(n-S)!S!\prod_i d_i!}.
 \tag{TI3.1}
\]
If nonzero $f,g$ have respective orders $e,e'$ in $w_\rho$, then
\[
 \deg\Gamma_{f,g}\leq S+o-\max(e,e'),\qquad
 [T^{S+o-e}]\Gamma_{f,g}
 =\frac{e!\binom{z+e}e B(\operatorname{in}_ef,\operatorname{in}_eg)}
 {o!S!\prod_i d_i!}\quad(e=e').
 \tag{TI3.2}
\]
Here $\operatorname{in}_e$ retains precisely coordinate degree $e$.

For a short proof, $[F,E]=r+\mathrm{Euler}_w+\mathrm{Euler}_y$ gives
the rational harmonic decomposition.  A finite binomial identity shows
that $\Phi$ acts on $E^l\ker F$ by $\binom{z+l}l/\binom Al$; hence
\[
 \Phi_{o,n,z}E^e=
 \frac{(z+1)^{\overline e}}{A^{\underline e}}
 E^e\Phi_{o-e,n-e,z+e}.
 \tag{TI3.3}
\]
Write $f=w_\rho^eh$, $q=\mathcal F(h)$ and $v=\mathcal F(g)$.
With $m=n-S$, the Leibniz formula
\[
 F^j(y_\rho^mq)=\sum_{u=0}^j\binom ju m^{\underline u}
 y_\rho^{m-u}\partial_{w_\rho}^uF^{j-u}q
\]
shows that $B(E^{e+j}F^j(y_\rho^mq),y_\rho^mv)/m!$ has degree at
most $S+j$.  Its leading coefficient is
$B(w_\rho^{e+j}\partial_{w_\rho}^jh,g)$.
The remaining coefficient in (TI3.3) is a polynomial of degree $o-e-j$.
At the top degree its sum reduces to
\[
 \sum_j\frac{(-w_\rho)^j}{j!}\partial_{w_\rho}^jh
 =h\big|_{w_\rho=0}.
\]
This proves both formulas (TI3.2), with symmetry supplying $\max(e,e')$.

\paragraph{TI4: integrality costs only logarithmically many powers of $p$.}
For integral $f,g$, let $L_0=\lfloor\log_p(S+o+r)\rfloor$ and
$L_3=\lfloor\log_p(3S)\rfloor$ when $S>0$. Then $p^K\Gamma$ has
integral binomial coefficients provided
\[
 K\geq(2r+1)L_0+(N+2)(r-1)L_3.
 \tag{TI4.1}
\]
To prove it, a multinomial coefficient with $r$ parts totaling $R$ has
valuation at most $(r-1)\lfloor\log_pR\rfloor$, by Legendre's formula.
Consequently $F^j\operatorname{fr}_n(f)/j!$ has coefficient valuations
at least $v_p(j!)-(r-1)L_0$.
Use self-adjointness in (TI3.1) and
$\binom Ao/\binom Aj=\binom{A-j}{o-j}/\binom oj$.
After accounting for the factorials in every row, the remaining factorial
valuation is exactly
\[
 v_p\binom{n-j}{S-j}-2v_p\binom oj-v_p\binom Sj.
\]
Evaluate only at $n=S,\ldots,2S+o\leq3S$.
This proves integrality with (TI4.1) at degree plus one consecutive integers;
Newton differences give integral coefficients in $\binom{T-S}j$, and
integer translation gives them in $\binom Tj$.
The case $S=0$ is constant.  This also proves the source's
$\texttt{GamInt}$ with $K=10(N+3)r^2(L_0+1)$.

\paragraph{TI5--TI6: why averaging removes most bad factors.}
Let $M(T)$ be an $n\times n$ matrix of integral binomial polynomials of
degree at most $d$, with nonzero determinant of degree $\delta$ and
leading coefficient of valuation $\lambda$.
On an integer progression $t_j=t_0+uj$, $p\nmid u$, $0\leq j<Y$, assume
$h\in\mathbb N$, $1\leq Y\leq p^h$ and $d\leq p^h$.
If $\beta(t)$ counts Smith factors of valuation at
least $E>0$, then
\[
 \lambda+\frac\delta{p-1}\leq nL
 \quad\Longrightarrow\quad
 \frac1Y\sum_j\beta(t_j)
 \leq n\left(\frac dY+\frac{L+dh/Y}{E}\right).
 \tag{TI5.1}
\]
Indeed $M(t+p^hX)$ is coefficientwise integral.  After Smith reduction at
$t$, its determinant coefficient of degree $k<\beta(t)$ has valuation
at least $E(\beta(t)-k)$.
Factoring the determinant over a splitting field gives
\[
 E\beta(t)\leq\lambda+\sum_\zeta\min\{v_p(t-\zeta),h\}
       +E\#\{\zeta:v_p(t-\zeta)\geq h\}.
\]
For each root the first sum averages to at most $1/(p-1)+h/Y$,
by counting residue classes modulo $p^i$; the small discs are disjoint,
so its last count sums to at most one. This proves (TI5.1).

For the actual matrix $p^{\lceil\kappa\rceil}\Gamma$, suppose the basis is
adapted to coordinate order and each order-$e$ initial Gram block, normalized
by $\prod_i d_i!$, has determinant valuation at most $\kappa_4$ times its rank.
Equal-order blocks alone supply the determinant's top coefficient in (TI3.2).
If $S,z\leq2a$, $r\geq2$, $\ell=\log_p(8a+r)+1$, and
$\kappa\geq(N+3)r^2\ell$, (TI4.1) supplies integral binomial coefficients.
Legendre's formula and (TI3.2) give
\[
 \frac1Y\sum_j\beta(n_j)
 \leq n\left\{\frac{2S}Y+
 \frac{\lceil\kappa\rceil+\kappa_4+3\ell+2S\ell/Y}{E}\right\},
 \tag{TI6.1}
\]
for $n_j=s+uj$, $Y,u\geq1$, $p\nmid u$, $S\leq s$ and $n_{Y-1}\leq2a$.
This keeps the integral normalization and the rounding term explicit.

\paragraph{TI7--TI9: the remaining bridge to the final exceptional function.}
The source next weights these estimates by varying highest-weight
multiplicities, supplies simultaneous integral partners, and constructs
paired maps into the cleaned rank-one edge lattices.
Those representation and lifting inputs remain open in this reconstruction.
Once their dimension-weighted inequality and numerical density bounds are
supplied, summing simultaneous degree translations gives
$\operatorname{mean}k_F\leq2\operatorname{mean}k_x$.
The actual final function has a lower-shift term (set to one at the boundary),
two upper-shift terms and $12Nkc/b'$.  Boundary and last-term costs at most
$1/Y$ each therefore give
\[
 \operatorname{mean}\kappa
 \leq6(C_A+C_s)NN_1/Y+2/Y
 \leq[6(C_A+C_s)+8]NN_1/Y.
 \tag{TI8.1}
\]
This final elementary transfer is proved in the companion.
It does not by itself establish the missing representation, block-density,
raw-lift or uniform-parameter hypotheses of the full Type-I conclusion.
